% TOP_PROBLEM: {"id": 3000023, "title": "Covering a symmetric crossing supermodular function with hyperedges of prescribed size", "category_id": 2, "category": {"id": 2, "name": "combinatorics", "display_name": "Combinatorics", "description": "Counting problems, graph theory, discrete structures.", "slug": "combinatorics", "order_index": 2, "created_at": "2026-07-31T15:26:25.671Z"}, "status": "open", "problem_number": "AMR-029-0023", "set_id": 13}
% FIRST_POSTED: 2026-10-01
% ENTRY_KIND: statement_only
% UnsolvedMath ID 3000023; source catalogue code AMR-029-0023
% !TeX program = lualatex
% Definition and sources only; this file is not an LLM solution attempt.
\documentclass[11pt,a4paper]{article}
\usepackage[margin=25mm]{geometry}
\usepackage{fontspec}
\setmainfont{lmroman10-regular.otf}[BoldFont=lmroman10-bold.otf,
 ItalicFont=lmroman10-italic.otf,BoldItalicFont=lmroman10-bolditalic.otf]
\usepackage{amsmath,amssymb,mathtools}
\usepackage{unicode-math}
\setmathfont{latinmodern-math.otf}
\AtBeginDocument{\let\setminus\smallsetminus}
\usepackage{xurl,hyperref}
\hypersetup{colorlinks=true,urlcolor=blue}
\setcounter{secnumdepth}{0}
\setlength{\parindent}{0pt}
\setlength{\parskip}{5pt}
\setlength{\emergencystretch}{3em}
\begin{document}
\section*{Covering a symmetric crossing supermodular function with hyperedges of prescribed size}\label{3000023}
{\small UnsolvedMath ID 3000023 · AMR-029-0023 · Combinatorics · AMR Open Problem Lists}
\subsection{Definitions and mathematical statement}
Let $V$ be finite and let $p:2^V\to\mathbb R$ be symmetric,
meaning $p(X)=p(V\setminus X)$. Assume crossing supermodularity:
$p(X)+p(Y)\le p(X\cap Y)+p(X\cup Y)$ whenever all four regions
determined by $X,Y$ are nonempty.
Specify positive integers $n_1,\ldots,n_k$.

Seek a criterion for the existence of an indexed hyperedge family
$H_1,\ldots,H_k\subseteq V$ satisfying
\[
 |H_i|=n_i\quad(1\le i\le k),\qquad
 |\{i:H_i\cap X\ne\varnothing,\ H_i\setminus X\ne\varnothing\}|
                                      \ge p(X)\quad(X\subseteq V).
\]
Each hyperedge meeting both sides of a cut contributes one unit,
irrespective of how many of its vertices lie on either side. Repeated
hyperedges are permitted and count with their indices.

This is a feasibility characterization problem, not the assertion that
every input is feasible. For example, $n_i>|V|$, a positive demand at
$\varnothing$, or an excessive cut demand makes realization impossible.
The result should provide sufficient as well as necessary conditions
on the prescribed sizes and demands. The sizes can differ; treating
only the equal-size case leaves the main additional difficulty.
The convention on cut contribution is essential: counting separated
vertex pairs of a hyperedge would define a different covering problem.
\subsection{Short English statement}
Characterize which symmetric crossing-supermodular cut demands
can be realized by a hypergraph with a prescribed list of hyperedge sizes.
\subsection{Sources}
\begin{itemize}
\item[S1] ULAM.AI and the UnsolvedMath Contributors, supplied problems.json, record 3000023 (AMR-029-0023), AMR Open Problem Lists. Catalogue identity and source transcription; CC BY 4.0.
\url{https://huggingface.co/datasets/ulamai/UnsolvedMath}.
\item[S2] Egres Open - List of open problems, Open-problem page: Covering a symmetric crossing supermodular function with hyperedges of prescribed size. Original source indexed by AMR-029-0023.
\url{https://oldlemon.cs.elte.hu/egres/open/List_of_open_problems}.
\item[S3] EGRES Open, Covering a symmetric crossing supermodular function with hyperedges of prescribed size. Individual source page and explanatory remarks.
\url{https://lemon.cs.elte.hu/egres/open/Covering_a_symmetric_crossing_supermodular_function_with_hyperedges_of_prescribed_size}.
\end{itemize}
\end{document}
